\begin{frame}{왜 하필 로봇인가: ``표''로는 안 되는 이유}
  표 기반 Q-테이블은 \textbf{상태를 행, 행동을 열}로 놓음.
  \vskip0.4em
  \kb{CartPole}의 상태: $(x_{\text{cart}},\,\dot{x}_{\text{cart}},\,
  \theta_{\text{pole}},\,\dot{\theta}_{\text{pole}})$ ---
  \textbf{4차원 연속 벡터}.
  \begin{itemize}
    \item ``$\cdots 0.011$ rad/s''인지 ``$\cdots 0.010$ rad/s''인지에 따라
      행이 \textit{별개}.
    \item 0.01 단위 그라데이션만 잡아도 차원당 $10^4$개,
      4차원이면 \textbf{$10^{16}$개의 행} ---
      \kb{차원의 저주}(curse of dimensionality)의 구체적 예.
  \end{itemize}
  \vskip0.4em
  로봇의 \textit{행동}도 연속 --- ``관절에 정확히 $+3.7$Nm''은
  이산 목록(열)에 없음. 연속값은 \textbf{무한히 많아서} 열 자체가
  존재하지 않음.
\end{frame}
